\begin{textsample}{POS Dim 4 – pb – Score 10.00 – pb/pb\_ef\_54.txt}  \label{ex:f4_pos_263}
9.1.2.
Direitos de aprendizagem e unidades temáticas para o Ensino \textbf{Religioso} No Ensino Fundamental, a BNCC normatiza que o Ensino \textbf{Religioso} adote a pesquisa e o diálogo como princípios mediadores e articuladores dos processos de observação, identificação, análise, apropriação e ressignificação de saberes, visando ao desenvolvimento dos seguintes direitos de aprendizagem : 1.
Conhecer os aspectos estruturantes das diferentes tradições/movimentos \textbf{religiosos} e filosofias de vida, a partir de pressupostos científicos, filosóficos, estéticos e éticos. 2.
Compreender, valorizar e respeitar as manifestações \textbf{religiosas} e filosofias de vida, suas experiências e saberes, em diferentes tempos, espaços e \textbf{territórios}. 3.
Reconhecer e cuidar de si, do outro, da coletividade e da natureza, enquanto expressão de valor da vida. 4.
Conviver com a diversidade de crenças, pensamentos, convicções, modos de ser e viver. 5.
Analisar as relações entre as tradições \textbf{religiosas} e os campos da cultura, da política, da economia, da saúde, da ciência, da tecnologia e do meio ambiente. 6.
Debater, problematizar e posicionar -se frente aos discursos e práticas de \textbf{intolerância}, \textbf{discriminação} e violência de cunho \textbf{religioso}, de modo a assegurar os direitos humanos no constante exercício da cidadania e da cultura de paz.
O Ensino \textbf{Religioso} está formatado na BNCC em três unidades temáticas.
A seguir, reproduzimos a lista dessas Unidades Temáticas e uma rápida explanação sobre os objetos de conhecimento aos quais se relacionam : 1. \textbf{Identidades} e Alteridades.
A unidade possibilita a distinção entre o “ eu ” e o “ outro ”, “ nós ” e “ eles ”, cujas relações dialógicas são mediadas por referenciais simbólicos : representações, saberes, crenças, convicções, valores a ser abordados especialmente nos anos iniciais do Ensino Fundamental.
Nessa unidade, pretende -se que os estudantes reconheçam, valorizem e acolham o caráter singular e diverso do ser humano, por meio da identificação e do respeito às semelhanças e diferenças entre o eu ( subjetividade ) e os outros ( alteridades ), da compreensão dos símbolos e significados e da relação entre imanência e transcendência. 2.
Manifestações \textbf{religiosas}.
Na unidade, pretende -se proporcionar o conhecimento, a valorização e o respeito às distintas experiências e manifestações \textbf{religiosas}, e a compreensão das relações estabelecidas entre as lideranças, denominações \textbf{religiosas} e as distintas esferas sociais, envolvendo símbolos, ritos, espaços, \textbf{territórios} sagrados e lideranças.
Nos símbolos, encontram -se dois sentidos distintos e complementares.
Por exemplo, objetivamente uma flor é apenas uma flor.
No entanto, é possível reconhecer nela um outro significado : a flor pode despertar emoções e trazer lembranças.
Assim, o símbolo é um elemento cotidiano ressignificado para representar algo além de seu sentido primeiro.
Sua função é fazer a mediação com outra realidade e, por isso, é uma das linguagens básicas da experiência \textbf{religiosa}.
Tal experiência é uma construção subjetiva alimentada por diferentes práticas ritualísticas, que incluem a realização de cerimônias, celebrações, orações, festividades, peregrinações, entre outras.
Enquanto linguagem gestual, os ritos narram, encenam, repetem e representam histórias e acontecimentos \textbf{religiosos}.
Desta forma, se o símbolo é uma coisa que significa outra, o rito é um gesto que também aponta para outra realidade.
Os rituais \textbf{religiosos} são geralmente realizados coletivamente em espaços e \textbf{territórios} sagrados ( montanhas, \textbf{mares}, \textbf{rios}, \textbf{florestas}, templos, santuários, caminhos, entre outros ), que se distinguem dos demais por seu caráter simbólico.
Estes espaços constituem -se em \textbf{lócus} de apropriação simbólico-cultural, onde os diferentes sujeitos se relacionam, constroem, desenvolvem e vivenciam suas \textbf{identidades} \textbf{religiosas}.
Nos \textbf{territórios} sagrados frequentemente atuam pessoas incumbidas da prestação de serviços \textbf{religiosos}.
Sacerdotes, líderes, funcionários, guias ou especialistas, entre outras designações, desempenham funções específicas : difusão das crenças e doutrinas, organização dos ritos, interpretação de textos e narrativas, transmissão de práticas, princípios e valores etc.
Portanto, os líderes exercem uma função pública e seus atos e orientações podem repercutir sobre outras esferas sociais, tais como economia, política, cultura, educação, saúde e meio ambiente. 3.
Crenças \textbf{Religiosas} e Filosofias de Vida.
Na unidade, são tratados os aspectos estruturantes das diferentes tradições/movimentos \textbf{religiosos} e filosofias de vida, particularmente sobre os mitos, ideia(s) de divindade(s), crenças e doutrinas \textbf{religiosas}, tradições orais e escritas, ideias de imortalidade, princípios e valores éticos.
Os mitos são outro elemento estruturante das tradições \textbf{religiosas}.
Eles representam a tentativa de explicar como e por que a vida, a natureza e o cosmos foram criados.
Apresentam histórias dos deuses ou heróis divinos, relatando, por meio de uma linguagem rica em simbolismo, acontecimentos nos quais as divindades agem ou se manifestam.
No enredo mítico, a criação é uma obra de divindades, seres, entes ou energias que transcendem a materialidade do mundo.
São representados de diversas maneiras, sob distintos nomes, formas, faces e sentidos, segundo cada grupo social ou tradição \textbf{religiosa}.
O mito, o rito, o símbolo e as divindades alicerçam as crenças, entendidas como um conjunto de ideias, conceitos e representações estruturantes de determinada tradição \textbf{religiosa}.
As crenças fornecem respostas teológicas aos enigmas da vida e da morte, que se manifestam nas práticas rituais e sociais sob a forma de orientações, leis e costumes.
Esse conjunto de elementos originam narrativas \textbf{religiosas} que, de modo mais ou menos organizado, são preservadas e passadas de geração em geração pela oralidade.
Desse modo, ao longo do tempo, cosmovisões, crenças, ideia(s) de divindade(s), histórias, narrativas e mitos sagrados constituíram tradições específicas, inicialmente orais.
Em algumas culturas, o conteúdo dessa tradição foi registrado sob a forma de textos escritos.
No processo de sistematização e transmissão dos textos sagrados, sejam eles orais, sejam eles escritos, certos grupos sociais acabaram por definir um conjunto de princípios e valores que configuraram doutrinas \textbf{religiosas}.
Estas reúnem afirmações, dogmas e verdades que procuram atribuir sentidos e finalidades à existência, bem como orientar as formas de relacionamento com a(s) divindade(s) e com a natureza.
No conjunto das crenças e doutrinas \textbf{religiosas}, encontram -se ideias de imortalidade ( \textbf{ancestralidade}, reencarnação, ressurreição, transmigração, entre outras ), que são norteadoras do sentido da vida dos seus seguidores.
Essas informações oferecem aos sujeitos referenciais tanto para a vida terrena quanto para o pós-morte, cuja finalidade é direcionar condutas individuais e sociais, por meio de códigos éticos e morais.
Tais códigos, em geral, definem o que é certo ou errado, permitido ou proibido.
Esses princípios éticos e morais atuam como balizadores de comportamento, tanto nos ritos como na vida social.
Também as filosofias de vida se ancoram em princípios cujas fontes não advêm do universo \textbf{religioso}.
Pessoas sem religião adotam princípios éticos e morais cuja origem decorre de fundamentos racionais, filosóficos, científicos, entre outros.
Esses princípios, geralmente, coincidem com o conjunto de valores seculares de mundo e de bem, tais como : o respeito à vida e à dignidade humana, o tratamento igualitário das pessoas, a liberdade de consciência, crença e convicções, e os direitos individuais e coletivos.
Cumpre destacar que os critérios de organização das habilidades na BNCC ( com a explicitação dos objetos de conhecimento aos quais se relacionam e do agrupamento desses objetos em unidades temáticas ) expressam um arranjo possível ( dentre

% matched lemmas: ancestralidade, discriminação, floresta, identidade, intolerância, lócus, mar, religioso, rio, território
\end{textsample}
